\centerline{\bf \Huge Турбо тренировка I}
\centerline{\bf \Huge }

\begin{flushright}
        \scriptsize{}
\end{flushright}

\problem{1}
Известно, что каждый посетитель дагестанской кофейни либо бородатый, либо безбородый. Поначалу процентная составляющая бородатых людей в этой кофейне была крайне мала, но  не равна нулю. После некоторого притока посетителей она стала свыше 99\%. Правда ли, что в какой-то момент она была ровно 99\%?

\problem{2}
Олег написал на доске десять натуральных чисел, среди которых нет двух равных. Известно, что из этих десяти чисел можно выбрать три числа, делящихся на 5. Также известно, что из написанных десяти чисел можно выбрать четыре числа, делящихся на 4 . Может ли сумма всех написанных на доске чисел быть меньше 75?

\problem{3}
Внутри равнобокой трапеции $A B C D$ с основаниями $B C$ и $A D$ расположена окружность $\omega$ с центром $I$, касающаяся отрезков $A B, C D$ и $D A$. Окружность, описанная около треугольника $B I C$, вторично пересекает сторону $A B$ в точке $E$. Докажите, что прямая $C E$ касается окружности $\omega$.

\problem{4}
Пусть $a,b > 0$, причём $a^5 - b^3\geqslant 2a$. Докажите, что $a^3 \geqslant 2b$